% Part: computability
% Chapter: recursive-functions
% Section: partial-functions

\documentclass[../../../include/open-logic-section]{subfiles}

\begin{document}

\olfileid{cmp}{rec}{par}
\olsection{توابع بازگشتی جزئی}

برای انگیزش تعریف توابع بازگشتی، توجه کنید که برهان ما دربارهٔ وجود
توابع محاسبه‌پذیری که بازگشتی اولیه نیستند در واقع چیزی بسیار بیشتر
را ثابت می‌کند. استدلال ساده بود: تنها از این واقعیت استفاده کردیم که
می‌توان توابع $f_0,f_1,\dots$ را چنان برشمرد که $f_x(y)$ به‌عنوان
تابعی از $x$ و~$y$ محاسبه‌پذیر باشد. پس استدلال دربارهٔ
\emph{هر رده از توابع که بتوان به این شیوه برشمرد} کاربرد دارد. این
ما را گرفتار می‌کند: می‌خواهیم توابع محاسبه‌پذیر را به‌صراحت توصیف
کنیم، اما هیچ توصیف صریحی از مجموعه‌ای از توابع محاسبه‌پذیر نمی‌تواند
جامع باشد!

راه برون‌رفت این است که بگذاریم توابع \emph{جزئی} وارد صحنه شوند.
خواهیم دید که برشمردن توابع محاسبه‌پذیر جزئی \emph{ممکن است}.
\iftag{TMs}{ در واقع، تقریباً از پیش می‌دانیم چنین است، زیرا می‌توان
  ماشین‌های تورینگ را به شیوه‌ای نظام‌مند برشمرد.}{} بعداً به استدلال
قطری خود بازمی‌گردیم و بررسی می‌کنیم چرا با گنجاندن توابع جزئی دیگر
پیش نمی‌رود.

اکنون پرسش این است: چه چیزی باید به توابع بازگشتی اولیه بیفزاییم تا
همهٔ توابع بازگشتی جزئی را به دست آوریم؟ باید دو کار انجام دهیم:
\begin{enumerate}
\item تعریف توابع بازگشتی اولیه را چنان تغییر دهیم که توابع جزئی را
  نیز مجاز بداند.
\item چیزی به تعریف \emph{بیفزاییم} تا چند تابع جزئی تازه نیز در آن
  گنجانده شوند.
\end{enumerate}

کار نخست آسان است. مانند پیش با صفر، جانشین و افکنش‌ها آغاز می‌کنیم
و نسبت به ترکیب و بازگشت اولیه بستار می‌گیریم. تنها تفاوت این است که
باید تعریف‌های ترکیب و بازگشت اولیه را چنان تغییر دهیم که امکان
تعریف‌نشده‌بودن برخی ترم‌های تعریف را مجاز بدانند. اگر $f$ و~$g$
توابع جزئی باشند، $f(x) \fdefined$ را برای این به کار می‌بریم که $f$
در~$x$ تعریف شده است، یعنی $x$ در دامنهٔ $f$ است؛ و
$f(x) \fundefined$ را برای خلاف آن، یعنی تعریف‌نشده‌بودن $f$ در~$x$،
می‌نویسیم. $f(x)\simeq g(x)$ یعنی یا $f(x)$ و~$g(x)$ هر دو
تعریف‌نشده‌اند، یا هر دو تعریف شده و برابرند. این نمادها را برای
ترم‌های پیچیده‌تر نیز به کار می‌بریم. قرارداد می‌کنیم که اگر $h$ و
$g_0$، \dots،~$g_k$ همگی توابع جزئی باشند، آنگاه
\[
h(g_0(\vec x),\dots,g_k(\vec x))
\]
تعریف شده است اگر و تنها اگر هر $g_i$ در~$\vec x$ تعریف شده باشد و
$h$ در $g_0(\vec x)$، \dots،~$g_k(\vec x)$ تعریف شده باشد. با این
تعبیر، تعریف ترکیب و بازگشت اولیه برای توابع جزئی همان تعریف بالاست،
جز اینکه باید «$=$» را با «$\simeq$» جایگزین کنیم.

آنچه برای به‌دست‌آوردن توابع جزئی به تعریف توابع بازگشتی اولیه
می‌افزاییم \emph{عملگر جست‌وجوی نامحدود} است. اگر $f(x,\vec z)$ هر
تابع جزئی روی اعداد طبیعی باشد، $\mu x\;f(x,\vec z)$ را چنین تعریف
می‌کنیم:
\begin{quote}
  کوچک‌ترین $x$ای که $f(0,\vec z)$، $f(1,\vec z)$، \dots،
  $f(x,\vec z)$ همگی تعریف شده‌اند و $f(x,\vec z)=0$، اگر چنین
  $x$ای وجود داشته باشد؛
\end{quote}
و قرار می‌گذاریم که در غیر این صورت $\mu x\;f(x,\vec z)$ تعریف‌نشده
باشد. این بند $\mu x\;f(x,\vec z)$ را به‌طور یکتا تعریف می‌کند.

\begin{explain}
توجه کنید که تعریف ما هیچ اشاره‌ای به\iftag{TMs}{ ماشین‌های تورینگ یا}{}
الگوریتم‌ها یا هیچ مدل محاسباتی مشخصی ندارد. اما مانند ترکیب و بازگشت
اولیه، شهودی عملیاتی و محاسباتی پشت جست‌وجوی نامحدود است. برای
محاسبه‌پذیری تابع جزئی، آرگومان‌هایی که تابع در آن‌ها تعریف نشده است
متناظر با ورودی‌هایی‌اند که محاسبه برایشان متوقف نمی‌شود. روش محاسبهٔ
$\mu x\;f(x,\vec z)$ چنین است: $f(0,\vec z)$، $f(1,\vec z)$،
$f(2,\vec z)$ و جز آن را محاسبه کنید تا مقدار~$0$ بازگردانده شود.
بااین‌حال، اگر هر محاسبهٔ میانی متوقف نشود، محاسبهٔ
$\mu x\;f(x,\vec z)$ نیز متوقف نمی‌شود.
\end{explain}

اگر $R(x,\vec z)$ هر رابطه‌ای باشد، $\mu x\;R(x,\vec z)$ را برابر
$\mu x\;(1\tsub\Char{R}(x,\vec z))$ تعریف می‌کنیم. به بیان دیگر،
$\mu x\;R(x,\vec z)$ کوچک‌ترین مقدار $x$ را بازمی‌گرداند که
$R(x,\vec z)$ برقرار است. پس اگر $f(x,\vec z)$ تابعی تام باشد،
$\mu x\;f(x,\vec z)$ همان $\mu x\;(f(x,\vec z)=0)$ است. اما توجه
کنید که تعریف اصلی ما کلی‌تر است، زیرا امکان می‌دهد $f(x,\vec z)$
همه‌جا تعریف نشده باشد (درحالی‌که تابع مشخصهٔ رابطه همواره تام است).

\begin{defn}
مجموعهٔ \emph{توابع بازگشتی جزئی} کوچک‌ترین مجموعه از توابع جزئی با
آریتی‌های گوناگون، از اعداد طبیعی به اعداد طبیعی است که صفر، جانشین و
افکنش‌ها را دربر می‌گیرد و نسبت به ترکیب، بازگشت اولیه و جست‌وجوی
نامحدود بسته است.
\end{defn}

البته برخی از توابع بازگشتی جزئی اتفاقاً تام‌اند، یعنی برای هر
آرگومان تعریف شده‌اند.

\begin{defn}
\ollabel{defn:recursive-fn}
مجموعهٔ \emph{توابع بازگشتی} مجموعهٔ آن دسته از توابع بازگشتی جزئی
است که تام‌اند.
\end{defn}

تابع بازگشتی را گاه برای تأکید بر اینکه همه‌جا تعریف شده است
«بازگشتی تام» می‌نامند.

\end{document}
